% OWNER: Arizona. Share the diagnostic/sensitivity page allocation.
\subsubsection{Accent association}
\label{sec:further-diagnostic-method}

A failure mode specific to one system could distort the accent comparison if
it varied across groups. Complementing the decoding checks, the following
analyses examine whether the literal \unk{} marker is linked to accent,
input modality or source wording (Table~\ref{tab:further-diagnostics}).
Checks that re-ran the model used the
Direct checkpoint, revision and decoding controls of Section~\ref{sec:systems}
on a Colab T4 GPU in FP16; exact library versions were not recorded. The
case-insensitive pattern \verb|<\s*unk\s*>| selects the same 508 archived
clips as exact-string matching. In all regressions, sentence length counts
whitespace-delimited words and sentence clusters use the original English
strings.

Logistic regressions of marker presence on accent, with and
without sentence length, tested the six accent coefficients (United States
English as reference) jointly by Wald tests, once with speaker-clustered and
once with sentence-clustered standard errors.

\subsubsection{Input modality}
The same checkpoint translated each clip's reference transcript
and Whisper transcript in batches of 32
(\texttt{SeamlessM4Tv2ForTextToText}). A speech re-run of 700
clips, randomly sampled with 100 per accent group (seed 760), kept the
generated token IDs to count special UNK ID~1, which decoding removes.

\subsubsection{Source wording}
Sentence-clustered logistic regressions related marker
presence to sentence length, proper-noun presence and the minimum
general-English Zipf frequency~\cite{speer2022wordfreq} of the clip's content
lemmas, the lower-cased dictionary forms of alphabetic tokens tagged NOUN,
VERB, ADJ or ADV (PROPN excluded) by spaCy
\texttt{en\_core\_web\_sm}~\cite{honnibal2020spacy}. The 48 clips without a
content token were assigned the median observed clip-level minimum (4.12).
Word removal used the 439 distinct English sentences of the 508 speech-marked
clips, restricted to the 385 whose text-input translation reproduced the
marker. At the decoding step where the first marker began, each whitespace-delimited
input word was deleted in turn, with the already generated output tokens
fixed. The selected word was the one whose deletion most reduced the summed
probability of the two alternative next tokens that begin the marker. Each sentence was
then translated afresh, without the fixed prefix, once without the selected
word and, as a control, once without another word chosen at random (seed
760). The outcome is the proportion of the 385 sentences whose new
translation contains no marker, with Wilson intervals. The Wilson intervals
treat clips in the special-ID check and sentences in word removal as
independent units.
